%% ma: L12-B14-0085
%% lop: 12
%% chuong: 5
%% bai: 14
%% dang: 12.14.12
%% loai: DS
%% muc: [VD, VD, NB, TH]
%% dap_an: "ĐĐĐS"
%% nguon: "Ảnh giáo viên gửi 07/11/2026 (Đề thi thử THPT Hoàng Diệu & Sào Nam, Đà Nẵng - 2026), Câu 4"
%% kien_thuc: "Bài 14 (phương trình mặt phẳng qua hai điểm vuông góc với mặt phẳng, hai mặt phẳng song song, khoảng cách giữa hai mặt phẳng song song)"
%% trang_thai: chinh-thuc
%% ghi_chu: "Tài liệu gốc không kèm đáp án; đáp án do Claude giải và kiểm bằng sympy."
\begin{ex}
Trong không gian $Oxyz$, cho hai điểm $A(1;-2;1)$, $B(0;1;-3)$ và mặt phẳng $(P)\colon x-y-3=0$.
\choiceTF
{\True Mặt phẳng $(Q)$ song song với mặt phẳng $(P)$, $(Q)$ cách $(P)$ một khoảng bằng $2\sqrt2$ và cắt trục $Ox$ tại điểm có hoành độ dương thì có phương trình $(Q)\colon x-y-7=0$.}
{\True Mặt phẳng $(\alpha)$ đi qua $A$, $B$ và vuông góc với mặt phẳng $(P)$ có phương trình $2x+2y+z+1=0$.}
{\True Điểm $A$ thuộc mặt phẳng $(P)$.}
{Mặt phẳng $(P)$ vuông góc với trục $Oz$.}
\loigiai{a) $(Q)\colon x-y+D=0$ ($D\ne-3$). $d\left((P),(Q)\right)=\dfrac{|D+3|}{\sqrt2}=2\sqrt2\Leftrightarrow D=1$ hoặc $D=-7$. Với $D=1$, $(Q)$ cắt $Ox$ tại hoành độ $-1$; với $D=-7$, tại hoành độ $7>0$. Vậy $(Q)\colon x-y-7=0$. Đúng.

b) $\overrightarrow{AB}=(-1;3;-4)$, $\vec n_P=(1;-1;0)$, $\left[\overrightarrow{AB},\vec n_P\right]=(-4;-4;-2)=-2(2;2;1)$. $(\alpha)\colon 2(x-1)+2(y+2)+(z-1)=0\Leftrightarrow 2x+2y+z+1=0$. Đúng.

c) $1-(-2)-3=0$ nên $A\in(P)$. Đúng.

d) $\vec n_P=(1;-1;0)$ không cùng phương với $\vec k=(0;0;1)$ nên $(P)$ không vuông góc với $Oz$ (thực ra $\vec n_P\cdot\vec k=0$ nên $(P)$ song song hoặc chứa $Oz$). Sai.}
\end{ex}
